The Processing Layer encapsulates \textbf{model-based and computer-vision} work for AgriHome, plus \textbf{asynchronous edge capture orchestration}. It is optional in the sense that core monitoring and CRUD remain available without external models; when present, it enriches the product with species hints, tray-level detections, post-ingest disease analysis, and scheduled Pi captures aligned with the SRS's emphasis on plant health insight and automatic image capture.

\subsection{Species and leaf inference adapter subsystem}

This adapter accepts a \textbf{plant or leaf image} (or derived crop) and returns a structured classification: proposed species or cultivar labels, confidence scores, and auxiliary metadata for the Application Layer to store alongside user-confirmed fields.

\begin{figure}[H]
\centering
\begin{tikzpicture}[font=\scriptsize, box/.style={draw, rounded corners, align=center, minimum width=2.4cm, minimum height=0.72cm}, arr/.style={-{Stealth}, thick}]
\node[box, fill=blue!8] (app) {Application};
\node[box, right=1.2cm of app, fill=orange!14] (inf) {Species\\adapter};
\node[box, right=1.2cm of inf, fill=gray!10, dashed] (gpu) {Model\\runtime};
\draw[arr] (app) -- node[above] {image} (inf);
\draw[arr, dashed] (inf) -- (gpu);
\end{tikzpicture}
\caption{Species inference adapter: hosted model runtime may be external.}
\label{fig:proc-species}
\end{figure}

\subsubsection{Assumptions}
\begin{itemize}
  \item Network latency between application and adapter is acceptable for interactive flows or runs asynchronously.
  \item Model versioning is tracked so historical reports can be interpreted correctly.
\end{itemize}

\subsubsection{Responsibilities}
\begin{itemize}
  \item Accept a well-defined inference request contract (image bytes or URL, optional tray/plant IDs).
  \item Run the trained classifier or delegate to a managed inference endpoint.
  \item Return normalized labels and scores without persisting domain state itself.
\end{itemize}

\subsubsection{Subsystem interfaces}
\begin{table}[H]
\centering
\small
\begin{tabular}{|p{0.9cm}|p{3.4cm}|p{3.0cm}|p{3.5cm}|}
\hline
\textbf{ID} & \textbf{Description} & \textbf{Inputs} & \textbf{Outputs} \\
\hline
PR-S1 & Application $\rightarrow$ Species adapter & Image + context & Label JSON \\
\hline
PR-S2 & Species adapter $\rightarrow$ Application & HTTP response & Parsed DTO \\
\hline
\end{tabular}
\caption{Species and leaf inference adapter interfaces}
\end{table}

\subsection{Tray vision and layout adapter subsystem}

The \textbf{tray vision} adapter analyzes a tray-level photograph to estimate plant count, regions of interest, or other layout features needed for monitoring dashboards. Outputs are merged into tray records or ephemeral UI state per product rules.

\begin{figure}[H]
\centering
\begin{tikzpicture}[font=\scriptsize, box/.style={draw, rounded corners, align=center, minimum width=2.4cm, minimum height=0.72cm}, arr/.style={-{Stealth}, thick}]
\node[box, fill=blue!8] (app) {Application};
\node[box, right=1.2cm of app, fill=orange!14] (tr) {Tray vision\\adapter};
\draw[arr] (app) -- node[above] {tray image} (tr);
\draw[arr] (tr) -- node[below] {detections} (app);
\end{tikzpicture}
\caption{Tray vision adapter: round-trip inference for tray imagery.}
\label{fig:proc-tray}
\end{figure}

\subsubsection{Assumptions}
\begin{itemize}
  \item Images meet minimum resolution and lighting heuristics expected by the model.
  \item A simulator or stub may substitute in development environments.
\end{itemize}

\subsubsection{Responsibilities}
\begin{itemize}
  \item Detect plants or regions; return geometry in a coordinate system agreed with the UI.
  \item Report confidence and failure modes explicitly.
\end{itemize}

\subsubsection{Subsystem interfaces}
\begin{table}[H]
\centering
\small
\begin{tabular}{|p{0.9cm}|p{3.4cm}|p{3.0cm}|p{3.5cm}|}
\hline
\textbf{ID} & \textbf{Description} & \textbf{Inputs} & \textbf{Outputs} \\
\hline
PR-T1 & Application $\rightarrow$ Tray adapter & Tray image & Detection set \\
\hline
\end{tabular}
\caption{Tray vision and layout adapter interfaces}
\end{table}

\subsection{Inference normalization subsystem}

\textbf{Normalization} converts heterogeneous model outputs into a single internal schema (units, label vocabularies, bounding-box ordering) so the Application Layer can persist results without adapter-specific branches scattered through domain code.

\begin{figure}[H]
\centering
\begin{tikzpicture}[font=\scriptsize, box/.style={draw, rounded corners, align=center, minimum width=2.5cm, minimum height=0.72cm}, arr/.style={-{Stealth}, thick}]
\node[box, fill=orange!12] (raw) {Raw model\\output};
\node[box, right=1.1cm of raw, fill=orange!8] (norm) {Normalizer};
\node[box, right=1.1cm of norm, fill=blue!8] (app) {Application};
\draw[arr] (raw) -- (norm);
\draw[arr] (norm) -- node[above] {canonical DTO} (app);
\end{tikzpicture}
\caption{Normalization: single canonical shape for downstream persistence.}
\label{fig:proc-norm}
\end{figure}

\subsubsection{Assumptions}
\begin{itemize}
  \item Each adapter documents its native schema version.
\end{itemize}

\subsubsection{Responsibilities}
\begin{itemize}
  \item Map labels to controlled vocabularies where required.
  \item Clamp or filter low-confidence results per policy.
  \item Provide stable fields for database columns and UI rendering.
\end{itemize}

\subsubsection{Subsystem interfaces}
\begin{table}[H]
\centering
\small
\begin{tabular}{|p{0.9cm}|p{3.4cm}|p{3.0cm}|p{3.5cm}|}
\hline
\textbf{ID} & \textbf{Description} & \textbf{Inputs} & \textbf{Outputs} \\
\hline
PR-N1 & Adapter $\rightarrow$ Normalizer & Vendor JSON & Canonical inference DTO \\
\hline
PR-N2 & Normalizer $\rightarrow$ Application & Canonical DTO & Ready for persist \\
\hline
\end{tabular}
\caption{Inference normalization subsystem interfaces}
\end{table}

\subsection{Edge disease hook subsystem}

The \textbf{edge disease hook} (\texttt{edge-disease-hook}) runs after a successful Pi ingest or direct capture when a plant ID is present and species inference is configured. It invokes the leaf/species adapter on the stored frame, writes prediction/report outcomes through domain services, and records monitoring events on skip or failure---without failing the ingest itself.

\begin{figure}[H]
\centering
\begin{tikzpicture}[font=\scriptsize, box/.style={draw, rounded corners, align=center, minimum width=2.4cm, minimum height=0.72cm}, arr/.style={-{Stealth}, thick}]
\node[box, fill=blue!8] (ingest) {Pi ingest\\complete};
\node[box, right=1.1cm of ingest, fill=orange!14] (hook) {edge-disease\\hook};
\node[box, right=1.1cm of hook, fill=orange!10] (cv) {Species\\adapter};
\draw[arr] (ingest) -- (hook);
\draw[arr] (hook) -- (cv);
\end{tikzpicture}
\caption{Edge disease hook: optional classification after Pi capture.}
\label{fig:proc-edge-disease}
\end{figure}

\subsubsection{Assumptions}
\begin{itemize}
  \item \texttt{CV\_SPECIES\_INFERENCE\_URL} (or equivalent) may be unset; then the hook records a skip event.
  \item Hook execution is fire-and-forget relative to the HTTP ingest response path when configured asynchronously.
\end{itemize}

\subsubsection{Responsibilities}
\begin{itemize}
  \item Load capture bytes and call plant analysis services.
  \item Persist prediction/report results when inference succeeds.
  \item Emit monitoring events for skip, success, and failure cases.
\end{itemize}

\subsubsection{Subsystem interfaces}
\begin{table}[H]
\centering
\small
\begin{tabular}{|p{0.9cm}|p{3.4cm}|p{3.0cm}|p{3.5cm}|}
\hline
\textbf{ID} & \textbf{Description} & \textbf{Inputs} & \textbf{Outputs} \\
\hline
PR-E1 & Ingest path $\rightarrow$ Edge disease hook & Capture + plant context & Monitoring / prediction side effects \\
\hline
PR-E2 & Edge disease hook $\rightarrow$ Species adapter & Image bytes & Diagnosis DTO \\
\hline
\end{tabular}
\caption{Edge disease hook subsystem interfaces}
\end{table}

\subsection{Capture schedule runner subsystem}

The \textbf{capture schedule runner} (\texttt{capture:schedule-runner}) is a periodic job that evaluates due \texttt{capture\_schedules}. For destination \texttt{raspberry-pi-edge}, it enqueues \texttt{capture\_now} commands (optionally with \texttt{runPoses: true}) onto linked devices so \texttt{agrihome\_agent} performs multi-angle capture without an operator present.

\begin{figure}[H]
\centering
\begin{tikzpicture}[font=\scriptsize, box/.style={draw, rounded corners, align=center, minimum width=2.4cm, minimum height=0.72cm}, arr/.style={-{Stealth}, thick}]
\node[box, fill=orange!12] (runner) {Schedule\\runner};
\node[box, right=1.1cm of runner, fill=gray!14] (cmd) {edge device\\commands};
\node[box, right=1.1cm of cmd, fill=teal!12] (agent) {agrihome\_agent};
\draw[arr] (runner) -- (cmd);
\draw[arr] (cmd) -- node[above] {heartbeat claim} (agent);
\end{tikzpicture}
\caption{Capture schedule runner: enqueues edge capture work for agents.}
\label{fig:proc-sched-runner}
\end{figure}

\subsubsection{Assumptions}
\begin{itemize}
  \item The runner is invoked on a cron-like cadence (for example every minute).
  \item Schedules and device linkage are authoritative in the Data Layer.
\end{itemize}

\subsubsection{Responsibilities}
\begin{itemize}
  \item Select due schedules with edge destinations.
  \item Insert pending capture commands with pose-run flags as configured.
  \item Avoid blocking the interactive Application Layer request path.
\end{itemize}

\subsubsection{Subsystem interfaces}
\begin{table}[H]
\centering
\small
\begin{tabular}{|p{0.9cm}|p{3.4cm}|p{3.0cm}|p{3.5cm}|}
\hline
\textbf{ID} & \textbf{Description} & \textbf{Inputs} & \textbf{Outputs} \\
\hline
PR-R1 & Runner $\rightarrow$ Data Layer & Due schedule query & Schedule/device rows \\
\hline
PR-R2 & Runner $\rightarrow$ Data Layer & \texttt{capture\_now} intents & Pending commands \\
\hline
\end{tabular}
\caption{Capture schedule runner subsystem interfaces}
\end{table}
